\section{Discussion}

The revised ARC-ADHD framework makes a narrower claim than the earlier integrated concept. It does not claim that adding more sensors, movement, AR, facial analysis, or unusual auditory stimulation should improve outcomes. Its central proposition is that repeated practice of a stable challenge-regulation-recovery rule can be measured within a digital intervention and can be tested for transfer beyond the trained task.

The strongest alternative explanation is engagement. A structured game may simply increase motivation or practice time without changing a transferable regulation process. Matching reward density, device exposure, contact time, and general visual quality in the active control is therefore essential. A second explanation is task familiarity: recovery may appear faster because the child learns the challenge sequence. Challenge-type variation, independent outcomes, and a confirmatory external endpoint are required to distinguish this from the proposed mechanism.

A third limitation is informant expectancy. Caregivers may observe the child using the intervention and may expect improvement. This can influence a caregiver-rated primary outcome even when the game has no specific regulation effect. Expectancy measurement, teacher report where feasible, and blinded independent assessment of selected outcomes should therefore be incorporated into later trials.

The framework is currently conceptual. No pilot or co-design data are presented. Statements about feasibility, acceptability, tolerability, and likely adherence are therefore hypotheses and development targets, not empirical findings. The one-author status of the current draft also limits the breadth of clinical and psychometric expertise represented. These limitations should be resolved through substantive collaboration before a participant-facing protocol is frozen.

The proposed mechanism also has an identification limit. A shorter \(T_{\mathrm{rec}}\) after repeated exposure would show change in the prespecified proximal process, and an association between recovery change and an independent outcome would support mechanistic coherence. Neither result alone proves causal mediation. A participant may improve on an external outcome through another pathway, or recovery telemetry may improve without transferable benefit. The statistical framework therefore treats proximal change, independent transfer, and a later confirmatory treatment effect as distinct empirical results rather than interchangeable evidence for one causal explanation.

Two design choices reduce avoidable measurement confounding without establishing construct validity by themselves. First, post-switch \(B\), \(S\), and \(V\) are normalized to a stable reference collected under the same target rule and GO-stimulus mapping rather than to the immediately preceding opposite rule. Second, \(T_{\mathrm{rec}}\) begins only when post-regulation recovery monitoring starts, while challenge-to-monitoring time \(T_{\mathrm{pre}}\) and total challenge-to-return time \(T_{\mathrm{total}}\) are retained separately. The remaining limitation is reference drift: an immutable within-session calibration reference may become less representative as practice or fatigue develops. Phase 0 must therefore quantify rule-specific reference stability across event order rather than silently refresh the baseline after seeing performance.

Bring-your-own-device delivery introduces a deliberate trade-off between ecological implementation and measurement control. Private-phone use can test whether the intervention is deployable under realistic hardware heterogeneity, but timing-sensitive measures may become less comparable across devices or after a device/profile change. Qualification, requalification, benchmark thresholds, coarse capability classes, and prespecified handling of timing discontinuities are consequently part of the measurement model rather than only software-engineering details.

Splitting the broader ARC-ADHD concept into companion studies is methodologically advantageous. Facial behavior has a distinct measurement-validity and privacy problem. Bilateral auditory stimulation has a distinct proof-of-mechanism problem. Embodied and AR interaction has distinct movement, spatial, device, and safety confounds. Wearable EEG has a separate acquisition, synchronization, artifact-control, and incremental-information problem. Treating these as separate causal or measurement questions makes null results interpretable and prevents feature accumulation from obscuring the primary mechanism.

The main unresolved measurement decision is the external primary endpoint. Candidate instruments differ materially in construct breadth, informant source, licensing, cultural adaptation, and overlap with ADHD symptoms. The endpoint should therefore be selected through a documented measurement audit rather than convenience.

Finally, \(T_{\mathrm{rec}}\) is useful only if its construction is reproducible. The feature set, reference window, stability criterion, and censoring rule cannot remain flexible after outcome inspection. The development programme explicitly treats those quantities as parameters that must be stabilized and frozen before inferential use. The accompanying supplementary methods specification extends this principle to the intervention manifest, challenge eligibility, controller replay, BYOD qualification, safety configuration, statistical analysis plan, and deviation record so that a later empirical report can be traced to a specific executable research state.
